% ====================================================================
% Section 1 (v2): Introduction.
% Written last, with the whole body in place. Voice: physics-first,
% honesty as a feature, equivalence-not-dominance.
% ====================================================================
\section{Introduction}
\label{sec:intro}

Network load balancing is a mature engineering discipline. Production systems
such as CONGA~\cite{conga} and DRILL~\cite{drill} spread traffic across a
network using carefully designed mechanisms: per-flowlet rerouting, multi-path
candidate enumeration, local queue sampling, and centrally or distributively
maintained congestion state. These systems work well, and improving on them is
hard.

This paper does not try to improve on them. It asks a different and, we think,
more interesting question:

\begin{quote}
\emph{How much of what engineered load balancers achieve can be recovered from
classical physics alone---from a packet treated as a particle obeying a couple
of elementary mechanical laws?}
\end{quote}

The motivation is partly methodological. An expert in routing reasons within
the conceptual vocabulary of routing; a physicist handed the same problem
reasons about gradients, forces, and conservation laws. Occasionally the second
framing reaches a familiar destination by an unfamiliar and more economical
road, and the comparison is informative regardless of which arrives first. Our
interest is therefore not a horse race but a question of \emph{parsimony}: how
little mechanism suffices to reach the operating point that elaborate
engineering reaches.

\subsection{The model, in one paragraph}

We treat a packet as a particle descending a potential field defined on the
network graph. The particle feels two forces. The first is a \emph{gradient}:
it rolls toward lower latency and lower congestion, subject to a hop budget that
lets it take a modestly longer, emptier detour rather than forcing every packet
onto the single shortest path. The second is a \emph{reaction} in the sense of
Newton's third law: when a packet is dropped on a congested edge, it pushes
back, depositing a decaying ``scar'' that steers subsequent packets away from
edges that have just failed. That is the entire router---two terms, no central
controller, no per-flow state, no precomputed path tables, sixty lines of code.

\subsection{What we find}

Our central result is deliberately stated as an \emph{equivalence}, not a
victory:

\begin{enumerate}
\item \textbf{Two terms suffice.} The two-term core is the survivor of an
ablation: an earlier model carried eight physically-motivated layers (effective
mass, a thermostat, a temporal ripple, buoyancy), and measurement showed that
only the gradient and the Newton-III scar earn their keep---the rest are inert
or mildly harmful (Section~\ref{sec:ablation}).

\item \textbf{It matches engineered routers.} Tested as a statistical
equivalence (TOST, margin $2$ percentage points) across fifteen topologies in a
flow-level simulator, the core is equivalent to DRILL on twelve of fifteen, and
equivalent to CONGA once CONGA is given a fair path budget. It reaches the same
operating point CONGA reaches only after enumerating and scoring sixteen to
thirty-two candidate paths per flow (Sections~\ref{sec:equivalence}--\ref{sec:fair}).

\item \textbf{A topological metric predicts when.} A quantity computable from
the graph alone, the tube/sp ratio (room to manoeuvre among near-shortest
paths), predicts where the equivalence holds and where engineered multi-path
routing keeps a genuine edge (Pearson $r = 0.78$); the two real WAN backbones,
with low tube/sp, are exactly the predicted exceptions
(Section~\ref{sec:applicability}).
\end{enumerate}

The upshot is that two classical forces, with none of the apparatus of modern
load balancers, reach the same place those balancers reach, on the topologies
where the geometry leaves room to do so---and that the qualifying clause is
itself predictable.

\subsection{A note on method}

This paper reports its negative results and its own corrections in full.
The eight-layer model is presented and then dismantled; an early, larger
advantage over CONGA is shown to have been an artifact of an unfair path budget
and is retracted (Section~\ref{sec:fair}). We regard this not as throat-clearing
but as the substance of the contribution: a claim that physics matches
engineering is worth precisely as much as the adversarial scrutiny it has
survived, and we would rather report an honest equivalence than an inflated
win. Every headline figure here has been checked against the strongest version
of its baseline we could build.
